\subsection{Bounded Flux as the Origin of Character Vanishing}
  \label{subsec:bounded-flux-character-vanishing}

  The bounded-flux condition introduced in Chapter~\ref{sec:spectral-admissibility}
  imposes a constraint on the maximal amplitude that a spectral mode of the relational
  Laplacian can sustain.
  This section clarifies why this constraint selects, at the level of representation
  theory, generators whose character vanishes in a common irreducible sector.

  Consider the relational Laplacian $\mathcal{L}_\chi$ acting on the configuration space
  of the $\chi$ substrate, and let $\rho$ denote an irreducible representation sector
  with associated Laplacian eigenvalue $\lambda_\rho$.
  The bounded-flux condition requires
  \begin{equation}
    A_\rho^{\max} \;=\; \frac{c_\chi}{\sqrt{\lambda_\rho}},
    \label{eq:admissibility-envelope-B}
  \end{equation}
  where $c_\chi$ is the maximal admissible projective speed.
  Sectors with smaller eigenvalues admit larger amplitudes and are therefore
  dynamically preferred.

  Now consider a generating set $S$ of a finite group $G$ acting on a spectral sector
  $\rho$.
  Each generator $s\in S$ acts on the eigenspace $V_\rho$ via the matrix $\rho(s)$.
  The local projective cost associated with $s$ is proportional to the real part of
  $1 - \chi_\rho(s)/\dim\rho$, where $\chi_\rho(s) = \mathrm{Tr}(\rho(s))$ is the
  character.

  The admissibility condition~\eqref{eq:admissibility-envelope-B} is saturated by the
  lowest available eigenvalue.
  For a Cayley graph with generating set $S$, the Laplacian eigenvalue in sector $\rho$
  takes the form
  \begin{equation}
    \lambda_\rho \;=\; |S| - \frac{1}{\dim\rho}\sum_{s\in S} \chi_\rho(s).
  \end{equation}
  Minimizing $\lambda_\rho$ subject to the constraint that $S$ generates $G$ is
  equivalent to maximizing $\sum_{s\in S}\chi_\rho(s)$.
  The global minimum is achieved when $\chi_\rho(s) = 0$ for all $s\in S$,
  in which case $\lambda_\rho = |S|$, and the admissibility envelope is at its most
  permissive.

  The character-vanishing condition $\chi_\rho(s)=0$ is therefore not an independently
  imposed algebraic requirement.
  It emerges as the necessary condition for generators to saturate the bounded-flux
  filter at the lowest possible spectral cost, and hence to belong to the admissible
  sector as defined by the projective structure of~$\Pi$.
